\documentclass[Screen16to9,17pt]{foils}
\usepackage{zencurity-slides}

\selectlanguage{english}

\usepackage{tabularx}

% ---------------------------------------------------------------------------
% Academy 2026 - Offensive Security subject
% Day 1 (second half): Network Attacks II -- hands-on & defence
% Maps to Subject plan objective 9 + education goal "Network resilience"
% Goals + Plan filled from the subject plan; content left as FIXME.
% ---------------------------------------------------------------------------

\begin{document}

\mytitlepage{Network Attacks II -- Hands-on \& Defence}
{Saint Cyber -- Cyber Akademiet, 2026}

\slide{Goals for today}
\begin{list2}
\item Demonstrate network-attack techniques in an authorised lab
\item Observe the effect of flooding on a target
\item Recommend defensive measures: patching, monitoring, rate/connection limits, filtering, CDN and DDoS protection
\end{list2}

\slide{Plan for today}
\begin{list2}
\item Lab environment and authorisation
\item hping3 -- crafting floods
\item Scapy -- packet crafting, task and challenge
\item Defences and mitigations
\item Red/blue exercise
\end{list2}

\slide{Time schedule}
% FIXME Keld: set timings for the half day
\begin{list2}
\item \textbf{FIXME:} 1) hping3 flooding -- XX min
\item \textbf{FIXME:} 2) Scapy + task -- XX min
\item \textbf{FIXME:} 3) Defences -- XX min
\item \textbf{FIXME:} 4) Red/blue exercise -- XX min
\end{list2}

\slide{Lab environment \& authorisation}
\begin{list2}
\item \textbf{FIXME:} describe the lab, target VMs, and the "only in the lab" rule
\end{list2}

\slide{hping3 -- crafting floods}
\begin{list2}
\item \textbf{FIXME:} example commands for SYN/UDP/ICMP floods -- test before class
\end{list2}

\slide{Observing the impact}
\begin{list2}
\item \textbf{FIXME:} Wireshark, nload, netstat / TCPview -- what to look for
\end{list2}

\slide{Scapy -- packet crafting}
\begin{list2}
\item \textbf{FIXME:} the basics the cadets need (IP()/TCP(), send/sendp, sniff)
\end{list2}

\slide{Scapy -- task}
\begin{list2}
\item \textbf{FIXME:} write the task (reuse your existing Scapy exercise)
\end{list2}

\slide{Scapy -- challenge}
\begin{list2}
\item \textbf{FIXME:} the spoofed-source / multi-port challenge
\end{list2}

\slide{Defences: patching \& monitoring}
\begin{list2}
\item \textbf{FIXME:} keep firmware/OS/software patched; detect early
\end{list2}

\slide{Defences: rate limits \& filtering}
\begin{list2}
\item \textbf{FIXME:} rate/connection limits, ingress filtering
\end{list2}

\slide{Defences: CDN \& DDoS protection}
\begin{list2}
\item \textbf{FIXME:} what a CDN does; scrubbing / DDoS protection services
\end{list2}

\slide{Red/blue exercise}
\begin{list2}
\item \textbf{FIXME:} the attack-and-defend exercise for the cadets
\end{list2}

\slide{Summary \& references}
\begin{list2}
\item \textbf{FIXME:} recap attacks + the matching defences
\item \textbf{FIXME:} add references / tool links
\end{list2}

\end{document}
